\section{Experiments}
\label{sec:experiments}

\subsection{Research Questions}

Three research questions assess our solution; a fourth is planned for Phase~3.

\paragraph{RQ1 (Effectiveness).} \emph{How accurately does our approach identify derived skill
variants, compared with exact hashing, similarity thresholds, and SkillReuse?} This question
tests whether lineage can be recovered from content at all, and whether our design improves on
existing detectors on the pairs that matter most: highly similar pairs that are not derived.
We report precision, recall, F1, and the Matthews correlation coefficient
(MCC)~\cite{chicco2020mcc} for the \emph{derived} class, computed on pooled out-of-fold
predictions, together with the standard deviation of F1 across folds. We also report the
false-positive rate on high-similarity pairs (non-derived pairs in bands B1--B3), precision and
recall per band, and, for our approach, a three-class confusion matrix and macro-F1.

\paragraph{RQ2 (Contributing signals).} \emph{Which signals contribute most to separating
derived variants from template siblings and unrelated skills?} We remove one feature group at
a time, retrain the best model, and report the change in F1 and in the high-similarity
false-positive rate. We also train a model on lexical features only.

\paragraph{RQ3 (Reuse concentration).} \emph{How concentrated is skill reuse when copies are
grouped into derivation families rather than by exact content hash?} For both units, exact
hash and family, we report the number of units, the share of units with a single file
occurrence, the Gini coefficient of occurrences per unit, the share of occurrences held by
the top 1\% of units, and the share of units and occurrences spanning two or more owners. We
also report the inflation ratio (distinct contents per family), the size of the largest
family, and the share of predicted derivations that commit history can orient in time.

\paragraph{RQ4 (Propagation, Phase 3).} \emph{When an original skill changes, do
its derived variants follow the change, stay frozen at an older version, or diverge?}

\subsection{Experimental Setup}
\label{sec:setup}

All experiments use the \ResGoldPairs{} gold pairs. We use five-fold stratified group
cross-validation; groups are formed by \ResGrouping{}, so no skill contributes to both
training and testing. We compare four baselines with our approach.

\emph{Normalised-text hash} predicts \emph{derived} only when the two bodies are identical
after lower-casing and collapsing whitespace, which is what exact hashing sees once trivial
formatting is ignored. \emph{Shingle Jaccard (MinHash)} thresholds the exact Jaccard similarity
of word 5-gram shingles, the quantity MinHash estimates. \emph{TF-IDF cosine} thresholds the
cosine similarity of TF-IDF vectors of the bodies. For both, the threshold that maximises F1
is chosen on the training folds only. \emph{SkillReuse-style} is our re-implementation of the
eight features of Zhu et al.~\cite{zhu2026skillreuse}: a flat TF-IDF similarity, LSA cosine
similarities for the YAML, prose, and code channels (set to zero when a channel is missing),
their three pairwise products, and a structural similarity over eight document statistics,
combined by a class-balanced logistic regression whose threshold is again tuned on the
training folds. It is trained on our gold set, so it is compared with our approach on equal
terms; it is not the authors' released implementation.

Our approach uses the 19 features of Table~\ref{tab:features} with three classifiers from
scikit-learn: logistic regression on standardised features, a random forest with 500 trees,
and histogram-based gradient boosting, all with balanced class weights and no threshold
tuning. The classifiers are trained on the three classes when the gold set has at least 15
template pairs, and on derived versus not derived otherwise.

\subsection{RQ1: Effectiveness}

\begin{table*}[t]
\caption{RQ1: detection of derived variants (pooled out-of-fold predictions; F1 $\pm$
standard deviation across folds). The last column is the share of non-derived pairs in bands
B1--B3 that a method labels as derived (lower is better).}
\label{tab:rq1}
\small
\restable{rq1_main.tex}{Run code/06\_rq1\_rq2.py to generate this table.}
\end{table*}

\begin{table*}[t]
\caption{RQ1 by similarity band: precision (P) and recall (R) of the strongest baseline and
of our best model.}
\label{tab:rq1bands}
\small
\restable{rq1_bands.tex}{Run code/06\_rq1\_rq2.py to generate this table.}
\end{table*}

\begin{table}[t]
\caption{RQ1: confusion matrix of our best model (out-of-fold).}
\label{tab:confusion}
\small
\restable{rq1_confusion.tex}{Run code/06\_rq1\_rq2.py to generate this table.}
\end{table}

Table~\ref{tab:rq1} compares all methods. The best configuration of our approach,
\ResOursName{}, reaches a precision of \ResOursP{}, a recall of \ResOursR{}, and an F1 of
\ResOursF{} (MCC \ResOursMCC{}) on the derived class. The strongest baseline,
\ResBaseName{}, reaches an F1 of \ResBaseF{}, and the SkillReuse-style detector \ResSRF{}. On
high-similarity pairs that annotators judged not derived, our approach labels \ResOursHiFP{}
as derived, against \ResBaseHiFP{} for the strongest baseline. Across the three classes it
reaches a macro-F1 of \ResMacroF{} (Table~\ref{tab:confusion}). Table~\ref{tab:rq1bands} shows
where the differences arise.
\tbd{Two or three sentences of error analysis from results/rq1\_oof\_predictions.csv: what do
the false positives and false negatives of the best model look like?}

\begin{quote}
\textbf{Answer to RQ1:} \tbd{One or two sentences summarising the result.}
\end{quote}

\subsection{RQ2: Contributing Signals}

\begin{table}[t]
\caption{RQ2: ablation of feature groups for our best model.}
\label{tab:rq2}
\small
\restablefit{rq2_ablation.tex}{Run code/06\_rq1\_rq2.py to generate this table.}
\end{table}

Table~\ref{tab:rq2} reports the ablation. Removing the \ResTopAblation{} group causes the
largest drop in F1 (\ResTopAblationDrop{}).
\tbd{Report what happens to the high-similarity false-positive rate when the boilerplate-aware
group is removed; this is the direct test of our design choice. Also compare with the
lexical-only model.}

\begin{quote}
\textbf{Answer to RQ2:} \tbd{One or two sentences.}
\end{quote}

\subsection{RQ3: Reuse Concentration}

\begin{table}[t]
\caption{RQ3: reuse concentration by exact hash and by derivation family. The last column
covers the whole in-scope population of the sample at the exact-hash level.}
\label{tab:rq3}
\small
\restablefit{rq3_table.tex}{Run code/07\_rq3.py to generate this table.}
\end{table}

Scoring all \ResCandPairs{} candidate pairs yields \ResDerivedEdges{} predicted derivations.
Grouping the \ResWorkHashes{} distinct contents of the working set into families leaves
\ResFamilies{} units, an inflation ratio of \ResInflation{}. \ResMultiFamilies{} families
contain two or more distinct contents, and the largest contains \ResLargestFamily{}, compared
with a largest connected component of \ResLargestCC{} before community detection. The Gini
coefficient of occurrences per unit moves from \ResGiniHash{} for exact hashes to
\ResGiniFam{} for families, and the share of units spanning two or more owners from
\ResCrossHash{} to \ResCrossFam{} (Table~\ref{tab:rq3}). Only \ResOrient{} of the predicted
derivations can be oriented in time from the dataset's commit history.

These figures are a lower bound. The official sample draws distinct contents uniformly, so most
members of a large family fall outside it and families appear smaller than they are. Phase~3
repeats the analysis on the full dataset.

\begin{quote}
\textbf{Answer to RQ3:} \tbd{One or two sentences.}
\end{quote}
